The geometric core of the argument is the generic pair-incidence fibre.
For a geometric generic point $x$ of
$\cM=(\cX/\cB)_{\mathrm{sm}}$, let $D_x$ be its unique effective divisor of
degree three.  The cyclic norm of a section of $H^0(C,3\kappa-D_x)$ is the
fixed norm line times a conic.  Its determinant has a structural linear
factor and a geometrically integral residual octic.  The factor-ordering
cover on that octic is split directly by the cubic norm: after restricting
to the line $Y=(-v/z)^3$, the residual conic is the product of the two
nontrivial cubic conjugates of the section.  This gives two rationally
labelled factors over $\CC$.  The scheme-theoretic Abel inverse converts
them into the two nonexceptional components of
\[
 C_x=\{y\in X:-x-y\in X\},
\]
and a separate boundary argument shows that no further nonexceptional
component occurs.

After spreading these two components over a nonempty open
$\cM^\circ\subset\cM$, top cohomology of the pair fibres and the
decomposition theorem produce an unshifted summand
$\IC_X\subset\Alt^2(\IC_X)$.  If $G=G(\IC_X)$ and
$V=\omega(\IC_X)$, the final input is therefore
\[
 \begin{gathered}
 G\subset \mathrm{O}(V),\qquad \dim V=7,\qquad V|_{G^0}\ \text{irreducible},\\
 \mu:\GGm\longrightarrow [G^0,G^0]
 \quad\text{with weights }-2,0,2\text{ of multiplicities }2,3,2,\\
 \operatorname{Hom}_G(V,\Lambda^2V)\ne0.
 \end{gathered}
\]
The exterior-square morphism rules out $\mathrm{SO}_7$, the Hodge cocharacter rules
out the seven-dimensional $\mathrm{PGL}_2$ module, and the same morphism excludes the
only possible disconnected extension of $G_2$.  Thus $G=G_2$.

\paragraph{Organization.}
The proof is organized by the following dependency chain.
\begin{enumerate}[label=\textup{(\arabic*)},leftmargin=*]
\item \cref{prop:theta-resolution,lem:surface-invariants,cor:theta-ic,lem:conormal-degree}
      establish the geometry of the theta section and $\chi(\IC_X)=7$.
\item \cref{prop:resolution-irregularity,prop:hodge-vector} compute the
      invariants of the resolution and the generic-twist Hodge
      multiplicities $(2,3,2)$.
\item \cref{prop:determinant-factorization,lem:residual-octic-integral,%
      prop:ordering-cover,prop:abel-incidence,lem:boundary-exhaustion}
      identify the two generic pair components, and
      \cref{thm:uniform-pair-splitting} spreads them.
\item \cref{lem:split-fibre-top,lem:off-theta-vanishing,lem:extract-ic}
      extract an unshifted $\IC_X$ summand in alternating pair convolution.
\item \cref{prop:connected-g2,cor:full-g2} identify the connected and full
      Tannaka groups.
\end{enumerate}

Unless a rational model or a finite-field reduction is explicitly mentioned,
the geometric construction is over $\CC$.  Intersection complexes,
cohomology, and Hodge modules retain their canonical $\QQ$-structures.  The
Tannakian category is obtained after extension of the underlying perverse
sheaves from $\QQ$ to $\CC$.  Thus all Tannaka groups and tautological
modules below are over $\CC$.
The convention $\PP(\mathcal E)=\operatorname{Proj}(\operatorname{Sym}\mathcal E^\vee)$ is used,
so $\PP(\mathcal E)$ parametrizes lines in the fibres of $\mathcal E$.  For
an integral scheme $Y$, $k(Y)$ denotes its function field and $k(y)$ the
residue field of a point $y$.

\section{Tannakian conventions}

Let $A$ be a complex abelian variety and let $m:A\times A\to A$ be addition.
For constructible complexes with coefficients in $\QQ$ or $\CC$, set
\[
K*L=Rm_*(K\bo L).
\]
Set
\[
 \begin{aligned}
 \mathsf S(A)&=\{Q\in\operatorname{Perv}(A,\CC):\chi(A,Q)=0\},\\
 \mathsf T(A)&=\{M\in D_c^b(A,\CC):{}^pH^i(M)\in\mathsf S(A)
                   \text{ for every }i\}.
 \end{aligned}
\]
The first is a Serre subcategory and the second is the corresponding
triangulated tensor ideal.  The quotient
\[
 \overline D(A)=D_c^b(A,\CC)/\mathsf T(A)
\]
inherits the perverse $t$-structure; its heart is
\[
 \overline{\operatorname{Perv}}(A)
   =\operatorname{Perv}(A,\CC)/\mathsf S(A).
\]
Convolution is $t$-exact on $\overline D(A)$, and the displayed heart is a
rigid symmetric monoidal neutral Tannakian category
\cite[Proposition~3.1]{JKLM2025}.  For the pure polarizable objects used
below, the tensor subcategory generated by one object is semisimple by the
decomposition theorem; its Tannaka group is therefore reductive.
Let $\delta_0$ denote the skyscraper perverse sheaf at the origin; its
image is the tensor unit for convolution.

For an integral subvariety $Z\subset A$ of dimension $d$, let
\[
 \IC_Z=j_{!*}\QQ_{Z_{\reg}}[d],
 \qquad j:Z_{\reg}\hookrightarrow A,
\]
be its perverse intersection complex, extended by zero to $A$.  If
$\chi(\IC_Z)>0$, let $G(\IC_Z)$ denote the Tannaka group of the category
generated by $\IC_Z\otimes_{\QQ}\CC$ in the displayed quotient, and by
$\omega(\IC_Z)$ its tautological module.  Then
\[
 \dim\omega(\IC_Z)=\chi(\IC_Z).
\]
The action on the tautological module is faithful by construction.  Scalar
extension is suppressed inside Tannakian expressions.  Let $\IC_Z^H$ denote
the polarizable pure Hodge module whose underlying
perverse $\QQ$-sheaf is $\IC_Z$.

The commutativity constraint contains the usual cohomological sign.  Since
the support in this paper has even dimension two, the alternating projector
on $\IC_X*\IC_X$ is carried by the fibre functor to the ordinary exterior
square of $\omega(\IC_X)$.  Let $\Alt^2(\IC_X)$ denote its image in the
Tannakian quotient.

\section{The shifted cyclic theta section}

Until \cref{sec:spreading}, fix $b=(A,B,c)\in\cB(\CC)$ and suppress it from
notation.  Thus
\[
 E:y^2=f_3(x),\qquad C:w^3=y-c,
\]
where $f_3(x)=x^3+Ax+B$.  Set
\[
 Y=y-c,
\]
so that $w^3=Y$.  The map $\pi:C\to E$ is a cyclic cover of degree three.
Its branch divisor is the reduced intersection of $E$ with the line $y=c$.

\subsection{The cover, the Prym, and the polarization}

Let $b_1,b_2,b_3\in E$ be the branch points and let $r_i\in C$ be the
ramification point above $b_i$.  Since
\[
 \operatorname{div}_E(y-c)=b_1+b_2+b_3-3O,
\]
the branch divisor belongs to $|3O|$ and, in the group law of $E$,
\begin{equation}
 b_1+b_2+b_3=O.
\label{eq:branch-sum}
\end{equation}
Put $\eta=\mathcal O_E(O)$.  The cyclic covering algebra is
\[
 \pi_*\mathcal O_C
 =\mathcal O_E\oplus\eta^{-1}\oplus\eta^{-2}.
\]
Riemann--Hurwitz gives $g(C)=4$, and the canonical bundle formula gives
\[
 K_C=\pi^*(K_E\otimes\eta^2)=\pi^*\eta^2=\kappa^2,
 \qquad \kappa=\pi^*\eta.
\]
By projection formula,
\[
 h^0(C,\kappa)
 =h^0(E,\eta)+h^0(E,\mathcal O_E)+h^0(E,\eta^{-1})=2.
\]

The pullback $\pi^*:J(E)\to J(C)$ is injective.  Indeed, if
$\alpha\in\operatorname{Pic}^0(E)$ and $\pi^*\alpha\simeq\mathcal O_C$, then
\[
 1=h^0(C,\pi^*\alpha)
  =h^0(E,\alpha)+h^0(E,\alpha\eta^{-1})
   +h^0(E,\alpha\eta^{-2}),
\]
and the last two terms vanish; hence $\alpha=\mathcal O_E$.  By duality,
$\ker(\Nm_\pi)$ is connected.  Thus set
\[
 P=\ker(\Nm_\pi).
\]
It has dimension three.  If $\tau(w)=\zeta_3w$, then
$\pi^*\Nm_\pi=1+\tau+\tau^2$ on $J(C)$.  Consequently
$\operatorname{Im}(1-\tau)\subset P$.  The $\tau$-invariant part of
$H^0(C,\Omega_C^1)$ is $\pi^*H^0(E,\Omega_E^1)$ and has dimension one, so
$\operatorname{Im}(1-\tau)$ has dimension three.  Hence
\begin{equation}
 P=\operatorname{Im}(1-\tau)=\ker(\Nm_\pi).
\label{eq:prym-conventions}
\end{equation}
The same argument works in the relative family.  Indeed
$\Nm_\pi\circ\pi^*=[3]$ on the relative elliptic Jacobian, so the norm map
is surjective; its kernel is smooth, and the fibrewise injectivity of
$\pi^*$ proved above makes every kernel fibre connected.  Thus
$\cP=\ker(\Nm_\pi)\to\cB$ is an abelian scheme of relative dimension
three.
